%% Chapter 4, Section 4.2: Proposal Kernels

\section{Proposal Kernels}
\label{sec:4.2}

We can classify Metropolis Hastings algorithms based on the form of their proposal kernels.
In this section, we study two types of proposal kernels that lead to the independence sampler
and the random walk Metropolis Hastings algorithm. Other proposal kernels will be studied in
subsequent chapters.

\subsection{Independence Sampler}
\label{sec:4.2.1}

Independence samplers are Metropolis Hastings algorithms where the proposal Markov kernel
does not include information on the current state of the chain. This is to say that, for some
distribution $g$ on $E$,
\[
  q_{\mathrm{Indep}}(x, z) = g(z),
\]
i.e.\ the proposal Markov kernel $q_{\mathrm{Indep}}(x, z)$ does not depend on the current state $x$ of the chain.

\begin{algorithm}[H]
\caption{Independence Sampler}
\label{alg:4.2}
\begin{algorithmic}[1]
\Require Target distribution $f$, initial distribution $\pi_0$, proposal $g(z)$, sample size $N$.
\State Define the Markov kernel
\[
  q_{\mathrm{Indep}}(x, z) := g(z).
\]
\State Run Metropolis Hastings (Algorithm~\ref{alg:4.1}) with inputs $f$, $\pi_0$, $q_{\mathrm{Indep}}(x,z)$, $N$.
\Ensure Sample $\{X^{(n)}\}_{n=1}^{N}$.
\end{algorithmic}
\end{algorithm}

Note that the acceptance probability of the independence sampler reduces to
\[
  a(x, z) = \min\!\left\{1,\, \frac{f(z)}{f(x)} \cdot \frac{g(x)}{g(z)}\right\}.
\]
By making analogy to importance sampling, we can define $w(x) := \frac{f(x)}{g(x)}$ and express the
acceptance probability as follows:
\[
  a(x, z) = \min\!\left\{1,\, \frac{w(z)}{w(x)}\right\}.
\]
Thus, moves $x \mapsto z$ for which $w(z) \geq w(x)$ are always accepted.

In practice, independence samplers typically perform poorly. However, they are an important
ingredient of the MCMC toolbox because their theoretical properties are well understood. To
gain intuition and to illustrate the flavor of the theory, we compare the independence sampler
with rejection sampling. Precisely, we establish a relation between the acceptance probability
for rejection sampling and the acceptance rate of the independence sampler at stationarity.

\begin{theorem}[Acceptance Rate of Independence Sampler]
\label{thm:4.4}
\textit{Suppose that the target $f$ and the
proposal distribution $g$ for an independence sampler satisfy, for some $M \geq 1$,}
\begin{equation}
  f(x) \leq M g(x), \quad \forall x \in E.
  \label{eq:4.1}
\end{equation}
\textit{Then, the independence sampler at stationarity has acceptance rate greater than $1/M$.}
\end{theorem}

\begin{proof}
Let $X \sim f$ and let $Z \sim g$. We work in the continuous setting so that $\Prob\bigl(f(z)g(x) = f(x)g(z)\bigr) = 0$. We need to show that
\[
  \Expect\bigl[a(X,Z)\bigr]
  = \Expect\!\left[\min\!\left\{1,\, \frac{f(Z)}{f(X)} \cdot \frac{g(X)}{g(Z)}\right\}\right]
  \geq \frac{1}{M}.
\]
We express the minimum with indicator functions
\[
  \min\!\left\{1,\, \frac{f(Z)}{f(X)} \cdot \frac{g(X)}{g(Z)}\right\}
  = \Ind{\frac{f(Z)}{f(X)} \frac{g(X)}{g(Z)} > 1}
  + \frac{f(Z)}{f(X)} \cdot \frac{g(X)}{g(Z)}\, \Ind{\frac{f(Z)}{f(X)} \frac{g(X)}{g(Z)} \leq 1}
\]
to obtain that
\begin{align*}
  \Expect\bigl[a(X,Z)\bigr]
  &= \iint \Ind{\{f(z)g(x) > f(x)g(z)\}}\, f(x)g(z)\, dx\, dz \\
  &\quad + \iint \frac{f(z)}{f(x)} \cdot \frac{g(x)}{g(z)}\, \Ind{\{f(z)g(x) \leq f(x)g(z)\}}\, f(x)g(z)\, dx\, dz.
\end{align*}
Now use symmetry and the bound~\eqref{eq:4.1} to deduce that
\begin{align*}
  \Expect\bigl[a(X,Z)\bigr]
  &= 2\iint \Ind{\{f(z)g(x) > f(x)g(z)\}}\, f(x)g(z)\, dx\, dz \\
  &\geq 2\iint \Ind{\{f(z)g(x) > f(x)g(z)\}}\, f(x)\,\frac{f(z)}{M}\, dx\, dz \\
  &= \frac{2}{M}\,\Prob\!\bigl(f(X_2)g(X_1) > f(X_1)g(X_2)\bigr) \\
  &= \frac{2}{M} \cdot \frac{1}{2} = \frac{1}{M},
\end{align*}
where $X_1, X_2 \iid f$.
\end{proof}

\medskip
\noindent\textbf{Pros and Cons.}
Theorem~\ref{thm:4.4} shows that the acceptance rate of the independence sampler at
stationarity is higher than the acceptance rate $1/M$ of rejection sampling. Moreover, contrary
to rejection sampling, the independence sampler does not throw away rejected draws. However, the independence sampler produces correlated draws and in practice is initialized far
from stationarity---several iterations will be needed to approximately reach stationarity (see
Theorem~\ref{thm:4.5} below). Among Metropolis Hastings algorithms, the independence sampler is
rarely used in practice and is mainly of theoretical importance.

As for rejection sampling, it is advisable to choose the independence sampler proposal $g$ to be
close to the target $f$, while at the same time being easy to sample from. Note that in the extreme
case where $g = f$, the algorithm outputs draws from $f$, the acceptance probability is 1, and the
Markov chain reaches stationarity immediately. However, choosing $g = f$ is clearly not useful
in the interesting case where sampling $f$ directly is not possible. The following result establishes
uniform ergodicity of the independence sampler. The notions of total variation distance and
uniform ergodicity are reviewed in Definitions A.23 and A.33.

\begin{theorem}[Uniform Ergodicity of Independence Sampler]
\label{thm:4.5}
\textit{Let $p_{\mathrm{Indep}}$ be the Markov kernel
of the independence sampler with proposal kernel $q_{\mathrm{Indep}}(x, z) = g(z)$. Suppose that there is
$M \geq 1$ such that $f(x) \leq Mg(x)$ for every $x \in E$. Then, for any $x \in E$,}
\[
  d_{\mathrm{TV}}\!\bigl(p_{\mathrm{Indep}}^n(x, \cdot),\, f\bigr)
  \leq \left(1 - \frac{1}{M}\right)^n.
\]
\end{theorem}

\begin{proof}
We only provide a sketch proof. If $x \neq z$,
\[
  p_{\mathrm{Indep}}(x, z)
  = q_{\mathrm{Indep}}(x, z)\, a(x, z)
  = g(z)\min\!\left\{1,\, \frac{f(z)}{f(x)} \cdot \frac{g(x)}{g(z)}\right\}
  = \min\!\left\{g(z),\, \frac{g(x) f(z)}{f(x)}\right\}
  \geq \frac{f(z)}{M}.
\]
The inequality $p_{\mathrm{Indep}}(x, z) \geq \frac{f(z)}{M}$ also holds when $x = z$. This bound allows us to conclude
the result with the same coupling argument used to show ergodicity of finite state space Markov
chains in Theorem A.25. A sketch is given below.

Let $\{B_n\}_{n=0}^\infty$ be a sequence of i.i.d.\ $\Bern\!\left(\frac{1}{M}\right)$ random variables, independent of all
other randomness, and define
\[
  W_{n+1} \sim
  \begin{cases}
    s(W_n, \cdot) & \text{if } B_n = 0,\\
    r(W_n, \cdot) & \text{if } B_n = 1,
  \end{cases}
\]
where $r(x, z) = f(z)$ and
\[
  s(x, z) = \frac{p_{\mathrm{Indep}}(x, z) - \frac{1}{M} f(z)}{1 - \frac{1}{M}}.
\]
The bound $p_{\mathrm{Indep}}(x,z) \geq \frac{f(z)}{M}$ ensures that $s$ defines a Markov kernel. Moreover, $\{W_n\}$ has
transition kernel $p_{\mathrm{Indep}}(x, z)$. The rest of the proof is identical to that of Theorem A.25.
\end{proof}

Notice the common underlying structure between Theorem~\ref{thm:4.5} and Theorem A.25. In both
cases we used a lower bound on the kernel $P$ in the sense that there exist $\delta > 0$ and a probability
measure $\nu$ such that
\[
  P^m(x, A) \geq \delta\,\nu(A), \quad \forall x \in E,\; \forall A \subset E.
\]
Indeed, this condition is equivalent to ergodicity of the chain. Whenever it holds,
\[
  d_{\mathrm{TV}}\!\bigl(P^n(x, \cdot),\, f\bigr) \leq (1-\delta)^{\lfloor n/m \rfloor}.
\]

\begin{remark}
\label{rem:4.6}
\textit{It can be shown that if $\mathrm{ess\,inf}\, \frac{g(z)}{f(z)} = 0$, then the independence sampler is not even
geometrically ergodic.}
\end{remark}

\medskip
\noindent\textbf{Scaling for the Independence Sampler.}
Suppose we want to use the independence sampler
with a proposal distribution chosen from some parametric family $\{g_\theta : \theta \in \Theta\}$. For example,
we may have decided to use a normal distribution as our proposal, and we want to determine an
appropriate mean and variance. This general type of problem is called ``scaling'' of Metropolis
Hastings methods.

Recall that we have shown that if $\frac{f(x)}{g(x)} \leq M$, then the independence sampler is uniformly
ergodic with rate $1 - \frac{1}{M}$. This suggests choosing the parameter $\theta$ which minimizes $M_\theta$ under
the constraint that $\frac{f(x)}{g_\theta(x)} \leq M_\theta$. Note that this is the same constrained optimization we would
like to (approximately) solve for in rejection sampling. In many problems of interest, finding an
upper bound for the ratio $f/g$ is unfeasible. In such cases, one can monitor the acceptance rate
of the independence sampler with various $\theta$'s, and choose the proposal $g_\theta$ that gives the highest
acceptance rate.

\subsection{Random Walk Metropolis Hastings}
\label{sec:4.2.2}

A Random Walk Metropolis Hastings (RWMH) algorithm is a particular type of Metropolis
Hastings method, where the proposal Markov kernel is chosen to be of the form
\[
  q_{\mathrm{RWMH}}(x, z) = g(z - x)
\]
for some distribution $g$.

\begin{algorithm}[H]
\caption{Random Walk Metropolis Hastings (RWMH)}
\label{alg:4.3}
\begin{algorithmic}[1]
\Require Target distribution $f$, initial distribution $\pi_0$, proposal distribution $g$, sample size $N$.
\State Define the Markov kernel
\[
  q_{\mathrm{RWMH}}(x, z) := g(z - x).
\]
\State Run Metropolis Hastings (Algorithm~\ref{alg:4.1}) with inputs $f$, $\pi_0$, $q_{\mathrm{RWMH}}(x,z)$, $N$.
\Ensure Sample $\{X^{(n)}\}_{n=1}^{N}$.
\end{algorithmic}
\end{algorithm}

The name RWMH comes from the fact that proposals are made according to a random walk:
proposed moves $Z^*$ may be obtained by sampling $\xi^* \sim g$ and setting
\[
  Z^* = X^{(n)} + \xi^*.
\]
The acceptance probability is
\[
  a(x, z) = \min\!\left\{1,\, \frac{f(z)}{f(x)} \cdot \frac{g(x-z)}{g(z-x)}\right\}.
\]
The original paper by Metropolis et al.\ proposed an RWMH with $g$ symmetric about $0$, which
simplifies the expression of the acceptance probability to
\[
  a(x, z) = \min\!\left\{1,\, \frac{f(z)}{f(x)}\right\}.
\]
Common choices of distribution $g$ are normal, Student's $t$-distribution, and uniform (in some
bounded subset, when $E$ is unbounded).

\begin{remark}
\label{rem:4.7}
\textit{If $E = \R^d$ and $q_{\mathrm{RWMH}}(x, z) = g(x-z) = g(z-x)$ for every $x, z \in E$, then the
RWMH kernel is not uniformly ergodic for any $f$. However, if the target $f$ is log-concave in the
tails and symmetric, and $g > 0$ is continuous, then the RWMH kernel is geometrically ergodic.
If $f$ is not symmetric, the conclusion still holds if additionally $g(x) \leq \ell e^{-\alpha|x|}$, where $\ell > 0$ and
$\alpha$ is some constant which quantifies the concavity of $\log(f)$.}
\end{remark}

\begin{figure}[htbp]
\centering
\includegraphics[width=0.95\textwidth]{figures/fig_04_1}
\caption{\textit{RWMH with standard Gaussian target $f = \Normal(0,1)$ and three proposals $g_\sigma = \Normal(0,\sigma^2)$ with $\sigma^2 = 0.024, 2.4, 240$. The choice $\sigma^2 = 2.4$ leads to good mixing of the chain, whereas $\sigma^2 = 0.024$ leads to low exploration and $\sigma^2 = 240$ leads to many rejections.}}
\label{fig:4.1}
\end{figure}

\medskip
\noindent\textbf{Scaling for RWMH.}
While for the independence sampler it is advantageous to maximize the
expected acceptance probability, this is \textit{not} the case for RWMH. Consider for intuition the case
where $f = \Normal(0,1)$ and $g_\sigma = \Normal(0, \sigma^2)$, as in Figure~\ref{fig:4.1}. Then,
\begin{itemize}
  \item Small $\sigma^2$ leads to high acceptance rate but poor exploration of the state space.
  \item Large $\sigma^2$ leads to good exploration but low acceptance rate.
\end{itemize}
There are some widely applicable guidelines for tuning proposals.
\begin{enumerate}
\item First, in the simple case of a normal target with a normal distributed random walk proposal,
the optimal choice of variance $\sigma^2$ can be determined by minimizing the integrated
autocorrelation time of the associated chains. The optimal value was found to be $\sigma^2 = 2.4$,
with corresponding acceptance rate
\[
  \alpha = \frac{2}{\pi}\arctan\!\left(\frac{2}{\sigma}\right) = 0.44.
\]
This result has motivated the general recommendation of aiming for an acceptance rate
around $1/2$ in low-dimensional problems.

\item Second, analysis of RWMH in the setting
\[
  f_d(x) = f(x_1) \times \cdots \times f(x_d), \quad x = (x_1, \ldots, x_d) \in \R^d,
  \qquad g_d(x) = \Normal(0, \sigma_d^2 I_d),
\]
reveals two important insights.

First, in the large-$d$ asymptotic, $\sigma_d^2$ should be scaled as $1/d$; this, in turn, implies that the
convergence time of the algorithm scales like $\calO(d)$.

Second, in a precise sense and under suitable assumptions, the optimal acceptance probability for large $d$ is roughly $0.234$. This result has motivated the general recommendation
of aiming for an acceptance rate around $1/4$ in high-dimensional problems.
\end{enumerate}
